\documentclass[a4paper,12pt]{article}
\usepackage[utf8]{inputenc}
\usepackage[T1]{fontenc}
\usepackage[ngerman]{babel}
\usepackage{amsmath, amssymb}
\usepackage{geometry}
\usepackage{parskip}
\usepackage{titlesec}
\usepackage{enumitem}
\usepackage{array}
\usepackage{booktabs}

\geometry{left=3cm, right=3cm, top=2.5cm, bottom=2.5cm}

\titleformat{\section}{\large\bfseries}{}{0em}{}[\titlerule]
\titleformat{\subsection}{\normalsize\bfseries\scshape}{}{0em}{}

% Glossar-Eintrag: Nummer, fetter Titel, Inhalt
\newcommand{\gentry}[3]{%
  \noindent\textbf{#1\quad #2}\nopagebreak\\[0.2em]
  #3\par\smallskip
}

\begin{document}

\begin{center}
  {\LARGE\bfseries Glossar zu Lektion 1}\\[0.5em]
  {\large Mathematische Grundlagen (Modul 61111)}
\end{center}

\vspace{1em}

Dieses Glossar fasst die wesentlichen Definitionen, Merkregeln und Sätze der
ersten Lektion kompakt zusammen. Nummerierungen entsprechen dem Lehrtext.

% ======================================================================
\section*{1.1 \quad Das Summensymbol $\Sigma$}

\gentry{1.1.1}{Merkregel (Vertauschen von Doppelsummen)}{%
  Wenn Summanden vertauscht werden dürfen (d.\,h.\ wenn Klammern um Summen
  weggelassen werden dürfen), so gilt:
  \[
    \sum_{1 \le i \le m} \sum_{1 \le j \le n} a_{ij}
    \;=\;
    \sum_{1 \le j \le n} \sum_{1 \le i \le m} a_{ij}.
  \]
  Die Summationsreihenfolge in Doppelsummen ist vertauschbar.
}

% ======================================================================
\section*{1.2 \quad Aussagen}

\gentry{}{Aussage}{%
  Ein Satz der Umgangssprache, dem eindeutig die Eigenschaft \emph{wahr} ($w$)
  oder \emph{falsch} ($f$) zugeordnet werden kann.
}

\gentry{1.2.3}{Negation $\lnot A$}{%
  Die Negation von $A$ ist die Aussage, die falsch ist, wenn $A$ wahr ist,
  und wahr, wenn $A$ falsch ist. Wahrheitstafel:
  \[
    \begin{array}{c|c}
      A & \lnot A \\\hline
      w & f \\ f & w
    \end{array}
  \]
}

\gentry{1.2.4}{Konjunktion $A \land B$}{%
  Genau dann wahr, wenn $A$ \emph{und} $B$ beide wahr sind.
  \[
    \begin{array}{cc|c}
      A & B & A \land B \\\hline
      w & w & w \\ w & f & f \\ f & w & f \\ f & f & f
    \end{array}
  \]
}

\gentry{1.2.5}{Disjunktion $A \lor B$}{%
  Genau dann wahr, wenn \emph{mindestens eine} der Aussagen $A$ oder $B$
  wahr ist.
  \[
    \begin{array}{cc|c}
      A & B & A \lor B \\\hline
      w & w & w \\ w & f & w \\ f & w & w \\ f & f & f
    \end{array}
  \]
}

\gentry{1.2.6}{Implikation $A \Rightarrow B$}{%
  Genau dann falsch, wenn $A$ wahr und $B$ falsch ist. $A$ heißt
  \emph{Prämisse}, $B$ heißt \emph{Konklusion}.
  \[
    \begin{array}{cc|c}
      A & B & A \Rightarrow B \\\hline
      w & w & w \\ w & f & f \\ f & w & w \\ f & f & w
    \end{array}
  \]
  Aus einer wahren Prämisse kann niemals etwas Falsches gefolgert werden.
}

\gentry{1.2.8}{Äquivalenz $A \Leftrightarrow B$}{%
  Genau dann wahr, wenn $A$ und $B$ denselben Wahrheitswert haben.
  \[
    \begin{array}{cc|c}
      A & B & A \Leftrightarrow B \\\hline
      w & w & w \\ w & f & f \\ f & w & f \\ f & f & w
    \end{array}
  \]
}

\gentry{1.2.9}{Wahrheitstafel}{%
  Eine Tabelle, die den Wahrheitswert einer zusammengesetzten Aussage in
  Abhängigkeit von den Wahrheitswerten der Teilaussagen auflistet.
}

\gentry{1.2.13}{Logische Äquivalenz}{%
  Zwei Aussagen heißen \textbf{logisch äquivalent}, wenn sie dieselben
  Wahrheitstafeln besitzen.
  \smallskip\\
  Wichtige logische Äquivalenzen:
  \begin{itemize}[topsep=2pt,itemsep=1pt]
    \item $A \Rightarrow B$ ist äquivalent zu $(\lnot B) \Rightarrow (\lnot A)$
          \quad (Kontraposition)
    \item $A \Leftrightarrow B$ ist äquivalent zu
          $(A \Rightarrow B) \land (B \Rightarrow A)$
    \item $\lnot(A \lor B)$ ist äquivalent zu $(\lnot A) \land (\lnot B)$
          \quad (De~Morgan)
    \item $\lnot(A \land B)$ ist äquivalent zu $(\lnot A) \lor (\lnot B)$
          \quad (De~Morgan)
  \end{itemize}
}

\gentry{1.2.15}{Existenzquantor $\exists$}{%
  Das Symbol $\exists$ wird durch „Es gibt ein" (im Sinne von „Es gibt
  mindestens ein") ausgedrückt. Schreibweise: $\exists\,x : A(x)$.
}

\gentry{1.2.17}{Allquantor $\forall$}{%
  Das Symbol $\forall$ wird durch „Für alle" ausgedrückt.
  Schreibweise: $\forall\,x : A(x)$.
}

\gentry{1.2.21}{Merkregel (Negation von Existenzaussagen)}{%
  Bei der Negation von $\exists\,x : A$ wird der Existenzquantor in einen
  Allquantor umgewandelt und die Aussage $A$ negiert:
  \[
    \lnot(\exists\,x : A) \;\equiv\; \forall\,x : \lnot A.
  \]
}

\gentry{1.2.23}{Merkregel (Negation von Allaussagen)}{%
  Bei der Negation von $\forall\,x : A$ wird der Allquantor in einen
  Existenzquantor umgewandelt und die Aussage $A$ negiert:
  \[
    \lnot(\forall\,x : A) \;\equiv\; \exists\,x : \lnot A.
  \]
}

\gentry{}{Vollständige Induktion}{%
  Sei $n_0 \in \mathbb{N}_0$ und $A(n)$ eine Aussage für alle $n \ge n_0$.
  Wenn gilt:
  \begin{enumerate}[label=(\roman*),topsep=2pt,itemsep=1pt]
    \item \textbf{Induktionsanfang:} $A(n_0)$ ist wahr.
    \item \textbf{Induktionsschritt:} Für alle $n \ge n_0$ gilt:
          $A(n)$ wahr $\Rightarrow$ $A(n+1)$ wahr.
  \end{enumerate}
  dann ist $A(n)$ für alle $n \ge n_0$ wahr.
}

\gentry{}{Beweisprinzipien}{%
  \begin{itemize}[topsep=2pt,itemsep=1pt]
    \item \textbf{Direkter Beweis:} Man nimmt $A$ an und schließt auf $B$.
    \item \textbf{Kontraposition:} Man beweist $\lnot B \Rightarrow \lnot A$
          statt $A \Rightarrow B$.
    \item \textbf{Widerspruch (indirekter Beweis):} Man nimmt $A$ und
          $\lnot B$ an und leitet einen Widerspruch her.
    \item \textbf{Ringschluss:} Um die Äquivalenz von (i), (ii), (iii) zu
          beweisen, genügt es, (i)$\Rightarrow$(ii), (ii)$\Rightarrow$(iii),
          (iii)$\Rightarrow$(i) zu zeigen.
  \end{itemize}
}

% ======================================================================
\section*{1.3 \quad Mengen}

\gentry{1.3.1}{Teilmenge, Gleichheit, leere Menge}{%
  \begin{itemize}[topsep=2pt,itemsep=1pt]
    \item $N \subseteq M$: jedes Element von $N$ liegt in $M$.
    \item $M = N$: $M \subseteq N$ und $N \subseteq M$.
    \item $\emptyset$: die Menge ohne Elemente.
  \end{itemize}
}

\gentry{1.3.3}{Mengenoperationen}{%
  Seien $M$, $N$ Mengen.
  \begin{itemize}[topsep=2pt,itemsep=1pt]
    \item \textbf{Vereinigung:} $M \cup N := \{x \mid x \in M \text{ oder } x \in N\}$
    \item \textbf{Durchschnitt:} $M \cap N := \{x \mid x \in M \text{ und } x \in N\}$
    \item \textbf{Differenz:} $M \setminus N := \{x \mid x \in M \text{ und } x \notin N\}$
    \item \textbf{Disjunkt:} $M$ und $N$ heißen disjunkt, falls $M \cap N = \emptyset$.
  \end{itemize}
}

\gentry{1.3.5}{Endliche und unendliche Mengen}{%
  Eine Menge $M$ heißt \textbf{endlich}, wenn sie nur endlich viele Elemente
  besitzt, andernfalls \textbf{unendlich}.
}

\gentry{1.3.7}{Mächtigkeit (Kardinalität)}{%
  \[
    |M| =
    \begin{cases}
      n, & \text{falls } M \text{ genau } n \text{ Elemente besitzt} \\
      \infty, & \text{falls } M \text{ unendlich ist.}
    \end{cases}
  \]
}

\gentry{1.3.9}{Produktmenge}{%
  Für nicht leere Mengen $M$, $N$:
  \[
    M \times N := \{(m,n) \mid m \in M,\, n \in N\}.
  \]
  $(m,n)$ heißt \textbf{geordnetes Paar}; $(m,n) = (m',n')$
  $\Leftrightarrow$ $m = m'$ und $n = n'$.
  Für endliche Mengen gilt $|M \times N| = |M| \cdot |N|$.
}

% ======================================================================
\section*{1.4 \quad Abbildungen}

\gentry{1.4.1}{Abbildung}{%
  Eine \textbf{Abbildung} $f: M \to N$ ist eine Teilmenge
  $f \subseteq M \times N$, sodass:
  \begin{enumerate}[label=(\roman*),topsep=2pt,itemsep=1pt]
    \item Für alle $m \in M$ gibt es ein $n \in N$ mit $(m,n) \in f$.
    \item Aus $(m,n) \in f$ und $(m,n') \in f$ folgt $n = n'$.
  \end{enumerate}
  $M$ heißt \textbf{Definitionsbereich}, $N$ heißt \textbf{Wertebereich}.
  Schreibweise: $n = f(m)$, \;$m \mapsto f(m)$.
}

\gentry{1.4.5}{Gleichheit von Abbildungen}{%
  $f: M \to N$ und $g: M' \to N'$ sind \textbf{gleich}, falls
  $M = M'$, $N = N'$ und $f(m) = g(m)$ für alle $m \in M$.
}

\gentry{1.4.6}{Bild und Urbild}{%
  Sei $f: M \to N$.
  \begin{itemize}[topsep=2pt,itemsep=1pt]
    \item $f(m)$: das \textbf{Bild} von $m$ unter $f$.
    \item $\mathrm{Bild}(f) = f(M) := \{n \in N \mid \exists\,m \in M: f(m) = n\}$:
          der \textbf{Bildbereich}.
    \item Ist $n \in \mathrm{Bild}(f)$, so heißt $m \in M$ mit $f(m) = n$
          ein \textbf{Urbild} von $n$.
  \end{itemize}
}

\gentry{1.4.9}{Injektiv, surjektiv, bijektiv}{%
  Sei $f: M \to N$.
  \begin{itemize}[topsep=2pt,itemsep=1pt]
    \item $f$ \textbf{surjektiv}: jedes $n \in N$ liegt im Bild von $f$.
    \item $f$ \textbf{injektiv}: jedes Element in $\mathrm{Bild}(f)$ besitzt
          genau ein Urbild
          (äquivalent: aus $f(m) = f(m')$ folgt $m = m'$).
    \item $f$ \textbf{bijektiv}: $f$ ist surjektiv und injektiv.
  \end{itemize}
}

\gentry{1.4.13}{Identische Abbildung}{%
  $\mathrm{id}_M: M \to M$, $m \mapsto m$. Sie ist bijektiv.
}

\gentry{1.4.14}{Komposition}{%
  Seien $f: L \to M$ und $g: M \to N$. Die \textbf{Komposition}
  $g \circ f: L \to N$ ist definiert durch $(g \circ f)(l) = g(f(l))$.
  \smallskip\\
  Die Komposition ist \textbf{assoziativ} (Prop.\,1.4.23):
  $(h \circ g) \circ f = h \circ (g \circ f)$.
  \smallskip\\
  Sind $f$ und $g$ bijektiv, so ist $g \circ f$ bijektiv.
}

\gentry{1.4.18}{Invertierbarkeit}{%
  $f: M \to N$ heißt \textbf{invertierbar}, wenn es eine Abbildung
  $f^{-1}: N \to M$ gibt mit $f^{-1} \circ f = \mathrm{id}_M$ und
  $f \circ f^{-1} = \mathrm{id}_N$.
}

\gentry{1.4.22}{Korollar (Charakterisierung invertierbarer Abbildungen)}{%
  Sei $f: M \to N$. Dann sind äquivalent:
  \begin{enumerate}[label=(\alph*),topsep=2pt,itemsep=1pt]
    \item $f$ ist bijektiv.
    \item $f$ ist invertierbar.
  \end{enumerate}
}

% ======================================================================
\section*{1.5 \quad Verknüpfungen}

\gentry{1.5.1}{Verknüpfung}{%
  Sei $M$ eine Menge. Eine \textbf{Verknüpfung} auf $M$ ist eine Abbildung
  $\circ: M \times M \to M$.
}

\gentry{1.5.5}{Eigenschaften von Verknüpfungen}{%
  Sei $\circ$ eine Verknüpfung auf $M$.
  \begin{itemize}[topsep=2pt,itemsep=1pt]
    \item \textbf{kommutativ}: $a \circ b = b \circ a$ für alle $a,b \in M$.
    \item \textbf{assoziativ}: $(a \circ b) \circ c = a \circ (b \circ c)$
          für alle $a,b,c \in M$.
    \item \textbf{neutrales Element} $e$: $a \circ e = e \circ a = a$
          für alle $a \in M$.
  \end{itemize}
}

\gentry{1.5.7}{Invertierbarkeit bzgl.\ einer Verknüpfung}{%
  Besitzt $\circ$ ein neutrales Element $e$, so heißt $m \in M$
  \textbf{invertierbar}, wenn es ein $m' \in M$ gibt mit
  $m \circ m' = m' \circ m = e$.
  Man nennt $m'$ \textbf{invers} zu $m$.
}

\gentry{1.5.10}{Distributivgesetze}{%
  Seien $\circ$ und $\#$ zwei Verknüpfungen auf $M$. Die
  \textbf{Distributivgesetze} lauten:
  \[
    m_1 \circ (m_2 \mathbin{\#} m_3) = (m_1 \circ m_2) \mathbin{\#} (m_1 \circ m_3),
    \qquad
    (m_1 \mathbin{\#} m_2) \circ m_3 = (m_1 \circ m_3) \mathbin{\#} (m_2 \circ m_3).
  \]
}

% ======================================================================
\section*{1.6 \quad Körper}

\gentry{1.6.1}{Körper}{%
  Ein \textbf{Körper} ist eine Menge $K$ mit zwei Verknüpfungen $+$ und $\cdot$
  sowie zwei verschiedenen ausgezeichneten Elementen $0$ und $1$, sodass gilt:
  \begin{enumerate}[label=(\roman*),topsep=2pt,itemsep=1pt]
    \item $+$ ist kommutativ, assoziativ; $0$ ist neutrales Element;
          jedes $a \in K$ besitzt ein additives Inverses $-a$.
    \item $\cdot$ ist kommutativ, assoziativ; $1$ ist neutrales Element;
          jedes $a \in K \setminus \{0\}$ besitzt ein multiplikatives
          Inverses $a^{-1}$.
    \item Distributivgesetz: $a \cdot (b+c) = a \cdot b + a \cdot c$.
  \end{enumerate}
  Beispiele: $\mathbb{Q}$, $\mathbb{R}$, $\mathbb{C}$, $\mathbb{F}_2 = \{0,1\}$.
  \smallskip\\
  $\mathbb{Z}$ ist \emph{kein} Körper (ganze Zahlen $\neq \pm 1$ nicht invertierbar bzgl.\ $\cdot$).
}

\gentry{1.6.4}{Proposition (Rechenregeln für Körper)}{%
  Sei $K$ ein Körper. Dann gilt:
  \begin{enumerate}[label=(\alph*),topsep=2pt,itemsep=1pt]
    \item $a \cdot 0 = 0$ für alle $a \in K$.
    \item Aus $a \cdot b = 0$ folgt $a = 0$ oder $b = 0$
          (\emph{Nullteilerfreiheit}).
    \item Das additive Inverse $-a$ ist eindeutig bestimmt.
    \item Für $a \neq 0$ ist das multiplikative Inverse $a^{-1}$
          eindeutig bestimmt.
  \end{enumerate}
}

\gentry{1.6.5}{Notation in Körpern}{%
  \begin{itemize}[topsep=2pt,itemsep=1pt]
    \item $-a$: additives Inverses zu $a$; $b - a := b + (-a)$.
    \item $a^{-1} = \frac{1}{a}$: multiplikatives Inverses zu $a \neq 0$;
          $\frac{b}{a} := b \cdot \frac{1}{a}$.
  \end{itemize}
}

\gentry{1.6.6}{Bemerkung}{%
  In jedem Körper $K$ gilt: $(-1) \cdot a = -a$ für alle $a \in K$.
}

% ======================================================================
\section*{2.1 -- 2.2 \quad Matrizen: Addition und Skalarmultiplikation}

\gentry{2.1.1}{Matrix}{%
  Eine \textbf{$m \times n$-Matrix} über einem Körper $K$ ist eine rechteckige
  Anordnung
  \[
    A = (a_{ij})_{\substack{1 \le i \le m \\ 1 \le j \le n}}, \qquad a_{ij} \in K.
  \]
  $a_{ij}$: Eintrag an der Stelle $(i,j)$.
  Für $m = n$: \textbf{quadratische Matrix}; $a_{ii}$: \textbf{Diagonalelement}.
  Bezeichnung: $M_{mn}(K)$.
}

\gentry{2.1.5}{Matrizenaddition}{%
  Für $A, B \in M_{mn}(K)$: $(A+B)_{ij} = a_{ij} + b_{ij}$.
  \smallskip\\
  $+$ ist kommutativ, assoziativ; neutrales Element ist die
  \textbf{Nullmatrix} $0 \in M_{mn}(K)$ (alle Einträge $= 0$);
  additives Inverses: $(-A) = (-a_{ij})$.
}

\gentry{2.2.1}{Skalarmultiplikation}{%
  Für $r \in K$, $A = (a_{ij}) \in M_{mn}(K)$: $rA = (r \cdot a_{ij})$.
  \smallskip\\
  Regeln (Prop.\,2.2.3): $(rs)A = r(sA)$;\; $(r+s)A = rA+sA$;\;
  $r(A+B) = rA+rB$;\; $1 \cdot A = A$.
}

% ======================================================================
\section*{2.3 \quad Matrizenmultiplikation}

\gentry{}{Matrizenmultiplikation}{%
  Seien $A \in M_{mn}(K)$ und $B \in M_{ns}(K)$. Das Produkt
  $AB \in M_{ms}(K)$ ist definiert durch:
  \[
    (AB)_{ik} = \sum_{j=1}^{n} a_{ij} \, b_{jk}.
  \]
  \textbf{Merkregel:} Zeile $i$ von $A$ auf Spalte $k$ von $B$ legen,
  komponentenweise multiplizieren und aufaddieren.
}

\gentry{2.3.5}{Einheitsmatrix}{%
  $I_m = (\delta_{ij}) \in M_{mm}(K)$ mit $\delta_{ij} = 1$ für $i = j$
  und $\delta_{ij} = 0$ für $i \neq j$.
}

\gentry{}{Regeln der Matrizenmultiplikation (Prop.\,2.3.6, 2.3.8)}{%
  \begin{itemize}[topsep=2pt,itemsep=1pt]
    \item Assoziativgesetz: $(AB)C = A(BC)$.
    \item $I_m A = A$, \; $A I_n = A$.
    \item Im Allgemeinen \emph{nicht} kommutativ: $AB \neq BA$.
    \item Distributivgesetze: $(A+A')B = AB+A'B$;\; $A(B+B') = AB+AB'$.
    \item $(rA)B = r(AB) = A(rB)$.
  \end{itemize}
}

\gentry{2.3.9}{Invertierbare Matrix}{%
  $A \in M_{nn}(K)$ heißt \textbf{invertierbar}, wenn es
  $A^{-1} \in M_{nn}(K)$ gibt mit $A^{-1}A = AA^{-1} = I_n$.
  \smallskip\\
  Das inverse Element ist eindeutig bestimmt (falls existent).
  Produkte invertierbarer Matrizen sind invertierbar:
  $(A_1 \cdots A_m)^{-1} = A_m^{-1} \cdots A_1^{-1}$.
}

% ======================================================================
\section*{3.1 -- 3.2 \quad Elementarmatrizen}

\gentry{3.1.1}{Elementare Zeilenumformungen}{%
  Sei $A \in M_{mn}(K)$.
  \begin{itemize}[topsep=2pt,itemsep=1pt]
    \item $Z_{ij}$: Vertausche Zeile $i$ und Zeile $j$ ($i \neq j$).
    \item $Z_i(r)$: Multipliziere Zeile $i$ mit $r \in K$, $r \neq 0$.
    \item $Z_{ij}(s)$: Addiere das $s$-fache von Zeile $j$ zu Zeile $i$
          ($s \in K$, $i \neq j$).
  \end{itemize}
}

\gentry{3.2.6}{Elementarmatrizen}{%
  Matrizen in $M_{mm}(K)$, die sich aus $I_m$ durch \emph{eine} elementare
  Zeilenumformung ergeben:
  \begin{itemize}[topsep=2pt,itemsep=1pt]
    \item $P_{ij} = I_m - E_{ii} + E_{ij} - E_{jj} + E_{ji}$
          \quad (entspricht $Z_{ij}$)
    \item $D_i(r) = I_m + (r-1)E_{ii}$
          \quad (entspricht $Z_i(r)$)
    \item $T_{ij}(s) = I_m + s E_{ij}$
          \quad (entspricht $Z_{ij}(s)$)
  \end{itemize}
  Dabei ist $E_{ij}$ die Matrix mit Eintrag $1$ an der Stelle $(i,j)$,
  sonst $0$.
  \smallskip\\
  \textbf{Prop.\,3.2.9:} Linksmultiplikation mit einer Elementarmatrix
  bewirkt die entsprechende Zeilenumformung an $A$.
}

\gentry{}{Prop.\,3.2.11 -- Invertierbarkeit von Elementarmatrizen}{%
  Elementarmatrizen sind invertierbar; ihre Inversen sind wieder
  Elementarmatrizen:
  \[
    P_{ij}^{-1} = P_{ij}, \qquad
    D_i(r)^{-1} = D_i\!\left(\tfrac{1}{r}\right), \qquad
    T_{ij}(s)^{-1} = T_{ij}(-s).
  \]
}

% ======================================================================
\section*{3.3 \quad Zeilenäquivalenz}

\gentry{3.3.2}{Zeilenäquivalenz}{%
  $A, B \in M_{mn}(K)$ heißen \textbf{zeilenäquivalent} ($A \sim_Z B$),
  wenn es Elementarmatrizen $E_1, \ldots, E_r \in M_{mm}(K)$ gibt mit
  \[
    A = E_1 E_2 \cdots E_r B.
  \]
  Eigenschaften: $A \sim_Z B \Leftrightarrow B \sim_Z A$
  (Symmetrie) und Transitivität.
}

\end{document}
